%%%%%%%%%%%%%%%%%%%%%%%%%%%%%%%%%
%
%       EXERCÍCIO
%
%%%%%%%%%%%%%%%%%%%%%%%%%%%%%%%%%

\definevalorquestao

\begin{exercicioBanco}[\valorquestao]
Sejam \(S_{1}\) e \(S_{2}\) os conjuntos soluções das seguintes equações modulares, respectivamente: 
\begin{multicols}{2}
    \begin{enumerate}[(i)]
        \item \(|x-3| = 2x-6\),
        \item \(|2x-4| = 4-2x\).
    \end{enumerate}
\end{multicols}
Determine \(S_{1} \cap S_{2}\).

% Define as alternativas
\newcommand{\alternativas}{%
    \begin{center}
        \begin{tabularx}{\textwidth}{XXX}
            (a) \(\{3\}\). &
            (b) \(\{2,3\}\). &
            (c) \(\varnothing\). \\[5pt]
            (d) \(\{x \in \bbR \mid x \le 2\}\). &
            (e) \(\bbR\).
        \end{tabularx}
    \end{center}
}

% Define a resposta correta
\newcommand{\resposta}{C}

% Lógica condicional para exibição
\ifdefstring{\atividade}{lista}{%
    \alternativas
    \vspace{0.5em}
    
    \noindent\textbf{Resposta:} letra \textbf{\resposta}.
}{%
    \ifdefstring{\modo}{objetivo}{%
        \alternativas
    }{%
        % modo = discursiva → não mostra alternativas
    }
}
\end{exercicioBanco}

